\documentclass[10pt, a4paper]{article}
% ===== Essential Packages =====
\usepackage[utf8]{inputenc}       % UTF-8 encoding
\usepackage[T1]{fontenc}          % Better font encoding
\usepackage{lmodern}              % Modern Latin Modern font
\usepackage{ctex}                 % Chinese support (must be loaded after fontenc)
\usepackage{subcaption}
% ===== Math & Symbols =====
\usepackage{amsmath}              % Advanced math
\usepackage{amssymb}              % Additional math symbols
\usepackage{mathrsfs}             % Ralph Smith's Formal Script
\usepackage{braket}               % Dirac bra-ket notation
\usepackage{pdfpages}
% ===== Graphics & Tables =====
\usepackage{graphicx}             % Image support
\usepackage{float}                % Better float control
\usepackage{wrapfig}              % Wrapped figures
\usepackage{tikz}                 % Vector graphics
\usepackage[table]{xcolor}        % Colored tables
\usepackage{array}                % Enhanced tables
\usepackage{multirow}             % Multirow tables
\usepackage{tabularx}             % Adjustable-width tables
\usepackage{booktabs}             % Professional table lines
\usepackage{siunitx}              % SI units in tables
\usepackage{caption}              % Better captions
\usepackage{adjustbox}            % Box resizing

% ===== Code Listings =====
\usepackage{listings}             % Code blocks
\usepackage{tcolorbox}            % Colored boxes
\tcbuselibrary{listings, breakable} % Listings in tcolorbox

% ===== Misc =====
\usepackage{lipsum}               % Dummy text
\usepackage{footmisc}             % Better footnotes
\usepackage[margin=1in]{geometry} % Page margins
\usepackage{diagbox}
\usepackage{ctex}
\author{张世航}
\title{SPR}
\begin{document}
	\maketitle
	\section{原始数据处理}
	\[\includegraphics[width=\textwidth, height=0.5\textheight, keepaspectratio]{SPR_curves_with_FWHM.pdf}\]
    根据图示\\
    
    \scalebox{0.7}{%
    	\begin{tabular}{|c|c|c|c|c|c|c|c|c|}
    		\hline
    		甘油浓度 & $0\%$ & $5\%$ & $10\%$ & $15\%$ & $20\%$ & $25\%$ & 未知1 & 未知2 \\
    		\hline
    		$(\lambda, amplitude)$ & $(644.3,76.619)$ & $(660.1,76.297)$ & $(671.9,74.774)$ & $(690.9,73.590)$ & $(709.0,70.730)$ & $(727.2,70.167)$ & $(666.1,70.925)$ & $(713.2,71.510)$ \\
    		\hline
    		折射率$n$&$ 1.3336$&$1.3401$&$1.3476$&$1.3564$&$1.3646$&$1.3721$&$1.3477$&$1.3661$\\
    		\hline
    		半高宽(FWHM)&$109.2nm$&$120.9nm$&$119.9nm$&$122.4nm$&$167.1nm$&$155.7nm$&$143.4nm$&$154.3nm$\\
    		\hline
    		半高(half-height)&$11.58nm$&$12.55nm$&$12.95$&$12.75nm$&$15.62$&$14.12nm$&$14.18nm$&$14.74nm$\\
    		\hline
    	\end{tabular}%
    }
    \section{计算未知样品}
    \subsection{利用峰值计算}
    利用不同浓度的数据计算两个未知的.\\
    \[\includegraphics[width=\textwidth, height=0.2\textheight, keepaspectratio]{2.pdf}\]
    \[k = 3.3154285714285736\]
    \[b = 642.4571428571429\]
    得到未知1为$7.131161668390196\%$,未知2为$21.337469837986884\%$的甘油.\\
    \subsection{利用折射率计算}
    我们利用不同浓度的折射率进行计算.\\
    \[\includegraphics[width=\textwidth, height=0.2\textheight, keepaspectratio]{4.pdf}\]
    我们带入两个未知样品的折射率
    \[\text{未知样品1浓度} = (635.6990839270819\times1.3477-847.2194411029856)\% = 9.51221430554267\%\]
    \[\text{未知样品2浓度} = (635.6990839270819\times1.3661-847.2194411029856)\% = 21.209077449801043\%\]
    \section{光纤表面等离子体共振(SPR)传感传感性能评价}
    \[\includegraphics[width=\textwidth, height=0.2\textheight, keepaspectratio]{3.pdf}\]\\
    对于$S = \frac{\Delta \lambda}{\Delta n}$,我们对各个已知样品进行最小二乘法拟合.
    得到$S = 2111.6406 nm$\\
    \scalebox{0.7}{%
    	\begin{tabular}{|c|c|c|c|c|c|c|c|c|}
    		\hline
    		甘油浓度 & $0\%$ & $5\%$ & $10\%$ & $15\%$ & $20\%$ & $25\%$ & 未知1 & 未知2 \\
    		\hline
    		$(\lambda, amplitude)$ & $(644.3,76.619)$ & $(660.1,76.297)$ & $(671.9,74.774)$ & $(690.9,73.590)$ & $(709.0,70.730)$ & $(727.2,70.167)$ & $(666.1,70.925)$ & $(713.2,71.510)$ \\
    		\hline
    		折射率$n$&$ 1.3336$&$1.3401$&$1.3476$&$1.3564$&$1.3646$&$1.3721$&$1.3477$&$1.3661$\\
    		\hline
    		半高宽(FWHM)&$109.2nm$&$120.9nm$&$119.9nm$&$122.4nm$&$167.1nm$&$155.7nm$&$143.4nm$&$154.3nm$\\
    		\hline
    		半高(half-height)&$11.58nm$&$12.55nm$&$12.95$&$12.75nm$&$15.62$&$14.12nm$&$14.18nm$&$14.74nm$\\
    		\hline
    		FOM&19.34 &17.47 &17.61&  17.25& 12.64 &13.56&14.73 & 13.69\\
    		\hline
    		
    	\end{tabular}%
    }\\我们可以得到\[\bar{FOM} = 15.784633872868678\]
\end{document}